\documentclass[11pt]{article}
\usepackage[a4paper,margin=2.3cm]{geometry}
\usepackage{amsmath,amssymb}
\usepackage{booktabs}
\usepackage{tabularx}
\usepackage{enumitem}
\usepackage[hidelinks]{hyperref}
\setlist{nosep,leftmargin=*}
\renewcommand{\arraystretch}{1.15}

\title{A Brain That Is Never Switched Off:\\
Binary Neurons, One Local Learning Rule, Innate Wiring and Speech Only Through Muscles\\[4pt]
\large Architecture, theory, and a test on an invented logic language}
\author{Zhiwen Li\thanks{ORCID 0009-0005-8289-6393. Code, reports and negative results:
\url{https://github.com/leeR1ven/concept-brain} (MIT licence; archived by Software Heritage). Correspondence via the
repository.}\\
Independent Researcher}
\date{November 2026 \quad (technical report, version 1)}

\begin{document}
\maketitle

\begin{abstract}
We describe \emph{concept-brain}, a research prototype of an artificial brain that lives one continuous life: it is
never reset, never frozen and never trained offline. It consists only of binary neurons wired at birth by one
generator (three-dimensional, local, random), a single local Hebbian learning rule with concave growth and one global
weight decay, and output only through simulated muscles. Sensory areas send one-to-one ``copies'' of their feature
neurons to a multi-layer prefrontal cortex (PFC) whose last layer predicts the next input; a thought is a first-layer
copy lit from inside; speaking is that prediction returned to the muscles. Reward is only a tool that strengthens or
weakens the links that produced an action in the last one to two seconds.
We give (i) the architecture as a short rule set; (ii) a theory showing that the Hebbian readout of the last layer is
the readout of an autoregressive language model under stated conditions, with signal-to-noise and timing analyses;
(iii) measurements of the brain built from the rule set; and (iv) a test on an invented thirteen-word logic language,
chosen because it is small, exact and reproducible. There the brain learns,
by only listening to others' dialogues and with no reward, to expect the right answer word after ``\texttt{T ?}'' within
20 dialogues and after ``\texttt{NOT T ?}'' within one or two hearings, once sentences are spoken at the pace of its
working-memory window; with hand-wired word--muscle links and an innate brake it answers in turn without talking over
others. We list what does not work yet, which is most of language.
\end{abstract}

\tableofcontents

\section{Introduction}

Large language models are trained offline by back-propagation and then frozen. A brain learns while it lives: every
moment it hears, thinks, moves and changes a little. The goal of this project, set by the author (``the designer'' in the
repository's notes), is an artificial intelligence that learns in real time in the same way, built from neurons only,
and that eventually talks with people and acts in the world. Its design rules are:

\begin{itemize}
\item \textbf{Pure neurons.} Binary neurons and connections only: no symbols, no lists, no lookup tables, no program
written for a particular task. Development-time substitutes are allowed only if an equivalent neural circuit is named
and later restored.
\item \textbf{One local learning rule and one global decay} for the whole brain; no back-propagation, no global
gradient.
\item \textbf{Always alive.} No reset between conversations, no sleep, no offline training phase.
\item \textbf{Output only through muscles.} Speaking is moving muscles; the model hears its own voice.
\item \textbf{Innate wiring first.} Every area is generated at birth by the same generator; learning works on top of it;
the first generation may receive hand-wired, registered connections (``generation 0''); later generations inherit
merged connections.
\end{itemize}

This report concerns thought and language. Our contributions are:

\begin{enumerate}
\item A rules-only specification of the architecture (Section~\ref{sec:arch}; the repository's
\texttt{docs/ARCHITECTURE.md} is the single authoritative version).
\item A theory (Section~\ref{sec:theory}): the last layer's Hebbian readout of ``what comes next'' equals the readout of
an autoregressive model when context codes overlap little; four premises that must hold; a signal-to-noise formula; a
proof that a prediction-error variant is the last-layer step of back-propagation; and timing results that later turned
out to decide what the model can learn.
\item Measurements of the brain built strictly from the rule set (Section~\ref{sec:built}): what a newborn holds, how
much activity stays bounded, where information is lost.
\item A test judged by function and by reading the weights, not by waiting for long runs (Section~\ref{sec:logic}):
an invented logic language in which the brain listens, thinks and talks.
\item A list of the designer's ideas with their current status (Appendix~\ref{app:ideas}) and a separate study on making
a back-propagation-trained language model learn continually with local rules (Appendix~\ref{app:llm2hebb}).
\end{enumerate}

All numbers in this report come from runs in the repository; each one is traceable to a script and a report file.
Earlier prototypes that used symbolic shortcuts are superseded by the rule set and are not reported here.

\section{Architecture}\label{sec:arch}

\subsection{Neurons and networks}
\begin{enumerate}
\item A neuron is either on or off. An inhibitory neuron is an ordinary neuron whose output weights are negative.
\item Codes are distributed: one thing lights many neurons; there is no ``one word, one neuron''.
\item Input layers use thermometer on/off pairs: exactly one of each pair is on, so an input layer is always half on;
``nothing there'' is also a reading.
\item \textbf{Every area is made by the same generator}: three-dimensional, local and random. Each neuron connects to a
few neurons of the layer before, near in the $y,z$ plane; birth weights are fixed. Each neuron's threshold is set from
the signal it receives, so similar inputs light similar neurons.
\item Deeper sensory layers are sparser. In the PFC, the last layer lights as many neurons as the input layer; the
middle layers are freer (about 30\% lit, Section~\ref{sec:core}).
\item The model's own voice and other people's voices light similar but different neurons; no label is needed.
\end{enumerate}

\subsection{Areas}
\paragraph{Sensory areas.} Hearing, vision and touch (including the force each muscle uses). In vision the centre has
more pixels than the edge; early visual layers also have local motion neurons wired at birth to the eye muscles. An
\emph{abstract concept area} receives the feature neurons of all sensory areas and learns with the same rule: things that
occur together or in the same context come to share codes, which is how invariance is learned (a person seen in the
centre and then drifting to the edge keeps the same concept neurons lit, which then link to the peripheral neurons).

\paragraph{Chain neurons.} Every feature neuron of every layer of every area sends a one-to-one copy to the PFC's first
layer and to the hippocampus; every outgoing chain has a returning chain, so a lit copy re-lights its source neuron.
Neurons near the input that are always on carry no feature and get no chain.

\paragraph{Prefrontal cortex.} A multi-layer recurrent network:
\begin{itemize}
\item the first layer is the copies of all areas; the middle layers are fixed random local networks made by the same
generator; the last layer's activity at one moment is linked by Hebbian learning to the first-layer neurons that come
next (``what comes next''); the last layer's internal links also learn (related units excite each other);
\item \textbf{depth:} the fewest layers that suffice, comparable to the three to four relay steps of a cortical column
\cite{douglas_martin_2004}; layers delay the output relative to the input and also combine features. Two layers are not
enough (the answer to ``\texttt{NOT T}'' is an exclusive-or of two features, Section~\ref{sec:logic});
\item a birth loop lets the last layer re-light the first layer and hold a thought;
\item \textbf{output is the first layer lit from inside.} Three sources light it: prediction, the loop, and hippocampal
recall (returning through the sensory areas). Lit copies re-light their sources through the return chains: in sensory
areas this is the inner voice or an imagined picture, in the motor area an intention. The strength of the return is a
knob between ``reason'' and ``instinct'';
\item excitability is context length; thinking is continuous and never restarts; speaking is caused by thinking, and
the outside world only influences thinking;
\item \textbf{working memory} has two parts: the delay through the layers (about five ticks at birth) and thoughts held by
the loop and by learned links (more than a dozen ticks after learning). A sentence must fit this window for the PFC to
link it to what follows (Section~\ref{sec:logic}); in the cortex, holding over seconds relies on recurrent
activity \cite{goldman_rakic_1995}.
\end{itemize}

\paragraph{Hippocampus.} Its own population, fed at birth by random connections from the copies, with no birth recurrence.
Its internal and recall links use the same Hebbian rule but grow fast and fade fast. A fixed number of neurons is lit per
tick. Recall re-lights neurons in the sensory areas, which then influence the PFC indirectly. It stores everything;
similar signals recall similar memories; long-term memory is learned by the cortex through repeated recall. No sleep, no
time ring.

\paragraph{Motor area and speech.} Output only through muscles. Each muscle has a group of force neurons (more lit, more
force), tonically active blocking neurons (``default: do not move'') and normally silent stop neurons whose chains let the
PFC learn ``finish'' and ``stop it''. Whether to speak is a competition with no hard threshold: when unsure, do not
speak. Sound and action are linked only through the PFC: the last-layer neurons that represent a sound drive the
first-layer copies of that sound's muscle forces, which drive the force neurons through the returns, against the
blockers. Two exceptions are innate: an input-layer $\to$ stop-neuron brake driven by the loudness of nearby sound (it can
only stop movement; it only gives a \emph{tendency} to take turns), and visual motion $\to$ eye muscles.

\paragraph{Amygdala (innate values).} Its own population, wired at birth to react to certain kinds and ranges of signals;
when lit it drives tendencies (approach, stop, attend) and is the source of reward signals. There may be several (pleasure,
fear).

\subsection{Learning}
\begin{itemize}
\item \textbf{Temporal Hebbian rule} \cite{hebb_1949}. When $B$ lights while $A$ has just been on or just turned off, $A\to B$ grows; an
$A$ that stays on while $B$ lights neither grows nor shrinks; $A$ re-lighting after $B$ weakens $A\to B$; $B$ turning
off changes nothing; neurons lit together do not strengthen each other.
\item \textbf{Each first-layer copy checks the links into it} (last layer $\to$ first layer). A copy that came from
outside but was not pushed by the last layer past threshold plus a margin of 2 strengthens the links from the last-layer
neurons that were just on; a copy pushed past threshold that did not come weakens them; otherwise nothing changes. The
change per tick (step 0.3) is shared by the lit neurons linked to it. In Breakout this raised the accuracy of predicting
the player's key from 52\% (learning only at onsets) to 87\% and then 92\%.
\item \textbf{Concave growth, no upper limit:} $\Delta w = \eta\, a_i a_j / (1 + w/w_\text{ref})^{p}$; weakening and decay
slow down by the same factor; about 1\% plasticity is always kept. Strong, much-used links are hard to change, so a
learned skill is not forgotten (metaplasticity \cite{abraham_bear_1996,yang_2009,fusi_2005,benna_fusi_2016}).
\item \textbf{One global decay:} learned links are multiplied by a factor at intervals and removed when too weak; birth
links do not decay; the hippocampus decays fastest and deeper layers slower.
\item \textbf{Time is order}: nothing represents ``how long ago''.
\item \textbf{Reward is only a tool.} Being right strengthens, being wrong weakens the links that produced an action
tendency in the last one to two seconds (last layer $\to$ action copies); other links are untouched, so the model still
remembers what was wrong but no longer does it. Each action's choice is the group lit at that moment (the press group
when held, the stop group when not; not jumping is a choice too); nearer choices change more; when upset, what pleasure
strengthened in the last seconds is also weakened. Good and bad come only from innate values, never from the amount of
response. The model learns mostly by listening; reward is a rare correction: reward more than punish, and punish clearly
only when it was sure and wrong. The credited choice must be the act itself (the choice made at the moment of answering).
\item \textbf{Structural plasticity} (slow): always-on neurons split, never-on neurons are removed; the criterion is
whether things can be told apart.
\item No special mechanism is written for any phenomenon.
\end{itemize}

\subsection{Life, world and offspring}
The model lives continuously; the young may make mistakes; the innate brake against talking over others stays. The world
keeps teaching: an adult pleased by ``almost right'' is the reward; there is no repeat-after-me training. Training order:
(1) the first generation receives hand-wired ``sound $\leftrightarrow$ the muscles that make it'' links in the PFC; (2) it
listens to many dialogues between others; (3) when it thinks of an answer, it says it. Every hand-wired group is
registered, and learnability is tested only on content that was never wired. \textbf{Evolution:} random birth $\to$ test
$\to$ keep the best $\to$ breed; each number of a child is 40\% the father's, 40\% the mother's and 20\% their mean, plus
a small random change; children are born with merged connections, still have to learn, but should learn faster; values
are innate and not merged. The whole brain is scaled together. The final product is pure neurons.

\subsection{Implementation details fixed by measurement}\label{sec:core}
\begin{itemize}
\item The PFC's middle layers are lit at about 30\%. When too sparse, each neuron is too selective and the same sentence
in another voice falls apart: sentence identity at the last layer (60 sentences, a fourth voice) is 0.06 at 2.6\%, 0.21
at 10\%, 0.41 at 30\% and 0.42 at 50\%.
\item The number of thoughts (copies lit from inside) is capped by the \emph{usual} external input count, not this tick's;
otherwise the budget is zero in quiet moments and a thought cannot light. Thoughts go only to copies not lit from outside.
Too many thoughts blur: a trained brain holds a word longest with a budget of 10\% of the usual input; at 100\% it holds
less.
\item The total push into a last-layer neuron from the last layer's own learned links is capped by the spread of the
input push across neurons at birth (about 1.6); otherwise after long practice the internal push is a thousand times the
input and the state is no longer set by what is heard.
\item Generation-0 links are stored at a ``well-used'' strength so daily learning does not wash them out; they keep the
temporal rule, not the self-check, so that thinking of a sound without saying it does not weaken them.
\item The self-check: strengthen until the push exceeds threshold by 2; weaken only when a copy was really pushed past
threshold and did not come.
\end{itemize}

\subsection{What is not used}
Back-propagation and global gradients; self-excitation that keeps a neuron on; refractory periods to drive thoughts;
``later inhibits earlier''; time rings and time cells; one word per neuron, IDs, lists, records; PMI gates, hub
inhibition, divisive gain, birth biases and similar patches; separate error-correcting readouts, basal ganglia or value
modules; reflexes and switches written for a test world; reward from the amount of response.

\section{Theory}\label{sec:theory}

\subsection{Formal model}
Every tick $t$ every neuron is 0 or 1. Let $x_t\in\{0,1\}^n$ be the copies of all feature neurons of all areas (taken
one-to-one from the areas' previous state). In the PFC
\begin{align*}
L1_t &= x_t \lor y_t, & y_t &= [\,d_t-u_t>\theta\,]\land \mathrm{top}_m(d_t-u_t), \quad d_t = G\,a_{t-1}+W\,a_{t-1},\\
L2_t &= [F_2 L1_{t-1}>\theta_2], & L3_t &= [F_3 L2_{t-1}>\theta_3],\\
a_t &= \mathrm{top}_{|L1_{t-1}|}\!\left(F_4 L3_{t-1} + g\,V a_{t-1} - u^4_t\right),
\end{align*}
where $y$ are the thoughts (copies lit from inside), $a$ the last layer, $G$ the fixed birth loop, $W$ the learned
last$\to$first readout, $V$ the last layer's learned internal links (capped: $g\sum_j V_{ij}\le\sigma_\text{in}$), $u$ each
unit's own usual level (a running mean of the push it receives), $m=\kappa\bar n$ the thought budget ($\kappa$ excitability,
$\bar n$ the usual external count), and $F_k$ fixed local random networks (each unit receives $K$ nearby inputs with weights
drawn from $U(0.5,1.5)$, thresholds set for a lit fraction $\rho_\text{mid}$). $W$ and $V$ learn with one local rule with
concave factor $\varphi(w)=1/(1+w)$; all learned links decay together at intervals. The hippocampus is
$s_t=\mathrm{top}_{k_h}(R\,H1_{t-1}+U s_{t-1}-u^s)$ with recall $h_t=\mathrm{top}_m(Q s_{t-1}-u^h)$ and $H1_t = x_t\lor h_t$.
Thoughts and recalls re-light their sources through the returns: inner voice, imagined pictures, intended movements.
Speaking is the PFC predicting the next moment's muscle forces.

\subsection{The PFC does what a language model does}\label{sec:equiv}
An autoregressive model compresses the past into a state $h_t=F(x_t,h_{t-1})$, reads out $Wh_t$ and feeds the most likely
next input back in. In this brain the last layer $a_t$ is the state ($F_4$ of the current input plus $V a_{t-1}$ carrying
the past), $W$ is the readout, and the most-pushed copy becomes the next input as a thought. Both are ``state
$\to$ predicted next input $\to$ use it as input''. If different contexts $c$ use non-overlapping sets $S_c$ of last-layer
units, the push into copy $j$ is $|S_c|$ times the count of ``$c$ followed by $j$'', so the most pushed copy is the most
frequent continuation of $c$, which is what a fully trained language model outputs. In matrix form
$W=\sum_t x_{t+1}h_t^\top$, which for orthogonal $h$ equals the least-squares readout $YH^\top(HH^\top)^{-1}$ up to scale,
the solution back-propagation converges to in the last layer. Random fixed middle layers suffice to separate contexts
(Cover's theorem \cite{cover_1965}; echo-state and liquid-state networks \cite{jaeger_2001,maass_2002}).

Four premises must hold:
\begin{enumerate}
\item different contexts overlap little in the last layer (the more they overlap, the more the readout says the most
common thing regardless of context);
\item link strengths follow frequencies instead of growing without bound (the global decay);
\item the most pushed copies really light as thoughts (the budget is set by the usual input count);
\item the last layer is set mainly by recent input (the echo-state property); the same $A$ must use roughly the same units
each time. Unbounded internal links break this: the internal push was 12 times the input push after 3,000 dialogues and
over 1,000 times after 79,000, hence the cap on $V$.
\end{enumerate}

\paragraph{Signal-to-noise.} For a pair $A\to B$ against a competitor $X$,
$\text{SNR} = (n(A,B)-n(A,X))(1-o)\,/\,\sqrt{(\text{var}_B+\text{var}_X)/K_p}$, with $o$ the overlap of different contexts in
the last layer and $K_p$ the number of lit last-layer units a copy receives. For the 20 most frequent pairs of a corpus:
a large newborn ($o=0.19$, $K_p=79$) gives a median SNR of 2.4 with 5 of 20 pairs reliable (SNR $>3$); a trained large brain
($o=0.10$, $K_p=98$) 3.9 and 14 of 20; a small newborn ($o=0.16$, $K_p=17$) 1.2 and 0 of 20. The SNR does not depend on how
long the brain learns, only on the code: interference comes from overlap.

\paragraph{Counting versus prediction error.} With a plain counting readout the push into a copy is a bigram smoother whose
kernel is the overlap of codes; when the kernel is flat the readout ignores the context. On 300 frequent syllables heard
18 times each by three voices, counting predicted the next syllable 11.6\% of the time, the same as ignoring the context
(11.1\%), however many times it listened. A readout with a local prediction error (strengthen what came but was not
predicted, weaken what was predicted but did not come) is exactly the stochastic-gradient step of the output layer, i.e.\
the last step of back-propagation; it converges to ridge regression and reached 31--32\% when the decay was strong. The
first-layer copies' self-check (Section~\ref{sec:arch}) is this rule.

\paragraph{Where identity is lost.} The ear's input bits identify 81\% of syllables by nearest class mean; after four fixed
random, increasingly sparse auditory layers only 38\% remain. Unsupervised winner-take-all prototypes with a local
``yield when winning too often'' term recover 19/25/30\% with 300/600/1,200 units. About ten directions separate syllables
more than a syllable varies with itself.

\paragraph{Timing.} With syllables lasting three ticks, the last layer identifies the syllable just heard 40--43\% of the
time four ticks after its onset, 49--55\% at five, 48--52\% at six and 28--35\% at seven: the signal needs four ticks to
reach the last layer. A link ``last layer one tick ago $\to$ copies that light now'' therefore cannot see the previous
syllable when the next one starts. Reading out the next syllable at the right moment gave 10.9\% on heard sentences in
order against 7.2\% shuffled (unheard: 9.2\% vs 4.9\%; an ideal reader that knows the syllable and counted pairs: 20.8\%).
This timing turned out to decide what the brain can learn (Section~\ref{sec:pace}).

\section{The brain built from the rules}\label{sec:built}

\subsection{Newborn}
A strict whole-brain build (one scale knob for every area, one-to-one chains, a count-preserving PFC with a closing loop,
competition-only motor) gives, at birth and before any learning:
\begin{itemize}
\item a heard question is still identifiable 12 silent ticks later: 32-way decoding 1.00 at the end, 0.64--0.71 after six
ticks and 0.34--0.40 after twelve (chance 0.03), with one non-learning local inhibitory neuron per patch and a
\emph{reciprocal} birth loop (each last-layer unit returns to the first-layer copies with the largest path weight into it)
of gain 1.0; random birth loops hold no better than 0.16 at twelve ticks;
\item counts are preserved layer to layer ($L_{k+1}/L_k = 0.90$--$1.05$, in silence and in speech); without local
inhibition and with a strong loop the PFC saturates (95\% on) or freezes;
\item activity inflation during learning comes only from the learned gains of the last layer's internal links and of the
hippocampus' internal links (fixed at 0.1);
\item in the speech brain (34 sounds), all PFC layers tell the sounds apart in another voice (layers 1--3: 1.00, last
layer: 0.97); a heard question is held 0.66 at its end, 0.31 after six ticks and 0.13 after twelve.
\end{itemize}

\subsection{Learning: what is kept and what is lost}
\begin{itemize}
\item \textbf{Hippocampus.} With a full step per co-occurrence, summed drive and fast forgetting (20,000 neurons, 50 lit),
a new sentence heard once pushes its next syllable 1.5--2 times more than the others and is gone after about three
dialogues: short-term memory. Without birth recurrence (fast weights only) the same input returns to the same code
(0.85--1.00), a question links to its answer across a pause (0.94--1.00) and an ordinary line heard once is completed
(0.84--0.90); episodes that start alike are not told apart (0.45--0.56), so one-shot facts of the form ``my X is called
\dots'' reach only 0.25--0.54.
\item \textbf{Long overhearing} (RTX 4090, 7 ms per tick, 5,000 dialogues in about 25 minutes): activity stays bounded
throughout. Clearing the hippocampus moves the PFC's rank of the right answer only from 11.5 to 12.5: what the PFC learned
is in its own links (cortical long-term memory). A hippocampal internal gain of 0.1 locks after 2,000 dialogues; 0.01 does
not.
\item \textbf{Regularity.} From the same brain state, the same sentence lights the same units (thoughts 0.61--0.72 Jaccard,
last layer 0.78--0.86; different sentences 0.15 and 0.25--0.28; another voice raises the thoughts' consistency from
0.17--0.23 at birth to 0.49--0.51 after 69,000 dialogues). After a gap, however, the last layer's state while listening
differs completely (0.00), while the input (0.50) and the middle layers (0.65--0.81) are stable. Setting the last layer's
internal links to zero restores it (0.53); with 15 or more quiet ticks before listening it is stable (0.58--0.69). The
internal links hold the content just heard for about ten ticks, and turns in dialogue are only four to ten ticks apart, so
each sentence is covered by the previous one.
\item \textbf{A large corpus.} 6.6 million Chinese dialogues (LCCC) converted to toned syllables and read by four
synthetic voices: after 1,000 overheard dialogues the PFC's next-syllable prediction is still at chance while a count model
reaches 14\%; activity is healthy. Possibly for the reason found in Section~\ref{sec:pace}: a sentence spoken at a pace that does not fit the PFC's
window cannot be linked to what follows; this has not been checked on the corpus.
\end{itemize}

\section{An invented logic language: listening, thinking and talking}\label{sec:logic}

\subsection{Setup}
Instead of a human language we use a small invented one: thirteen words (\texttt{T}, \texttt{F}; \texttt{NOT},
\texttt{AND}, \texttt{OR}, \texttt{SAME}, \texttt{IS}; the names \texttt{A}--\texttt{D}; ``\texttt{.}'' and ``\texttt{?}'')
plus two answer words \texttt{YES} and \texttt{NO}. Each word is a fixed random pattern of lit cells on a small sheet seen
by the vision area. Five stages add sentences while the earlier ones stay:

\begin{center}\small
\begin{tabular}{lll}
\toprule
stage & sentences & needs\\
\midrule
1 & \texttt{T ?}, \texttt{F ?} & truth words\\
2 & \texttt{NOT T ?} & negation\\
3 & \texttt{T AND F ?}, \texttt{F OR T ?} & two values\\
4 & \texttt{A SAME B ?} & hold the first name, compare\\
5 & \texttt{A IS T . A ?} & bind a name to a value, ask later\\
\bottomrule
\end{tabular}
\end{center}

The brain answers with two muscles, YES and NO. Each muscle has its own press and stop (inhibiting) groups, so ``none'',
``one'' and ``both'' are all outputs; only exactly the right muscle is right. Answers are balanced per kind of sentence
(otherwise ``\texttt{AND} $\to$ always NO'' is right three times in four). The first runs use a $12\times16$ sheet (2,311 copies),
the later ones a four-times larger brain ($24\times32$ sheet, 9,372 copies); the PFC has four layers unless stated.

\subsection{Reward alone, with parameters computed from the weights}
With the trace kept only for the choice at the moment of answering and full credit, one reward moves that input's push by
about $3.6\times$ the step, so a step of 0.5 moves it by twice the motor noise and one mistake should change the next answer.
With settings taken from an earlier game test (step 0.06, a fifteen-frame trace, credit by surprise) stage 1 was not learned in
1,000 sentences; with the computed settings it was learned in 40--100. The two-layer PFC learned stage 1 in 239--352
sentences and never stage 2; the four-layer PFC learned stage 1 in 73 and stage 2 in 762. Readability of the answer at the
moment of answering (ridge, newborn small brain, 1.0 perfect, 0.5 chance) shows why:

\begin{center}\small
\begin{tabular}{lccccc}
\toprule
PFC & stage 1 & 2 & 3 & 4 & 5\\
\midrule
two layers & 0.88 & 0.56 & -- & -- & --\\
three layers, loop $\times2$ & 0.99 & 1.00 & 0.73 & 0.82 & 0.81\\
four layers & 1.00 & 1.00 & 0.85 & 0.48 & 0.92\\
four layers, loop $\times2$ & 1.00 & 1.00 & 0.97 & 0.90 & 0.98\\
\bottomrule
\end{tabular}
\end{center}

Halving the motor noise and keeping the PFC's own links fixed while the reward learned brought stage 2 down from 1,981 to
138 sentences; a four-times larger brain with a doubled loop learned all five stages once (stage 5 at sentence 3,263, about
five minutes of compute). Mixed together, however, the answers stayed at 60--75\%. Without noise, reading the push per kind
of sentence showed shortcuts (``\texttt{AND} $\to$ NO'', ``\texttt{OR} $\to$ YES''; \texttt{SAME} and \texttt{IS} near chance),
although one readout of the trained last layer still separated all answers at 0.92--0.94; the reward links had grown to a
mean of 3 (maximum 15) and kept about 30\% of their plasticity.

\subsection{Listening, as the architecture prescribes}
Following the training order of Section~\ref{sec:arch}, the brain only listens to others' dialogues (``\texttt{T ? YES}'',
``\texttt{NOT T ? NO}'', \dots), muscles still, no reward, and the PFC learns ``what comes next'' as always. A test without
learning asks which answer word the PFC expects just before it appears (its push into the YES-word copies against the
NO-word copies, and which word's copies take the thought budget).
\begin{itemize}
\item Stage 1: the expectation is right 100\% after 20 dialogues (step 0.3; steps 1, 3 and 10 are not faster and less
stable).
\item \textbf{Talking.} Generation-0 links (registered): each copy of an answer word lit from inside (thought of, not
heard) pushes its muscle's press group; tonic blockers keep the muscles still by default; the innate brake pushes both
stop groups by the number of lit input cells, holding over the gap between two words. In its own turn after a question,
a listened brain thinks of the right word about 11 copies per tick against 4.4 for the wrong one out of a budget of 35; the
link strength was computed from these numbers. Half of life is others' dialogues, half is its own turns; its own word shows
on the sheet more softly (its own voice).
\end{itemize}
\begin{center}\small
\begin{tabular}{lcccc}
\toprule
stage 1, about 500 turns & talks over others & echoes others & answers in its turn & right when it answers\\
\midrule
newborn, no reward & 0\% & 0\% & 80\% & 75\%\\
after listening, no reward & 0\% & 0\% & 40\% & 99\%\\
after listening, rare correction & 7\% & 0\% & 83\% & 96\%\\
\bottomrule
\end{tabular}
\end{center}
Turn-taking appears as a tendency from fixed links, as intended. Without any correction over 10,000 episodes (5,000 of
them its own turns) the accuracy fell to 50--65\%: the brain hears its own wrong words and learns them as ``what comes next''. With rare correction (reward when
right, punish only when sure and wrong) it learned negation and \texttt{AND}/\texttt{OR} (stage 2 after 2,734--7,500 episodes,
stage 3 after 14,962; half of the episodes are its own turns), slowly, and eagerness to answer grew until it often answered the others' questions before they did.

\subsection{The pace must fit the working-memory window}\label{sec:pace}
By listening alone the brain first learned ``\texttt{NOT} $\Rightarrow$ answer the opposite'' and got the plain sentences
wrong. The cause was not the rule. Recording the last layer just before the answer word in a brain that does not learn, and
then letting the copies' self-check rule learn on these fixed states, isolates the question:
\begin{itemize}
\item said slowly (three frames per word plus a blank), \texttt{NOT} has left the last layer when the answer comes: the
states of ``\texttt{T ?}'' and ``\texttt{NOT T ?}'' are nearly the same units, and the four sentences were never all
learned in 60 passes, with excitatory or with signed links, whatever the loop (1--4$\times$), the excitability
(10--100\%) or the length of the quiet between dialogues;
\item said fast with no pause, \texttt{T} has not yet reached the last layer;
\item \textbf{one frame per word, no gap, a pause of three frames before the answer}: ``\texttt{T ?}'' lights 52 units and
``\texttt{NOT T ?}'' 44, of which 43 and 35 are their own. Same size, different units: the four sentences are learned in
\textbf{one to two passes}, with the unchanged rule and excitatory links only.
\end{itemize}
Listening at that pace (learning on), stages 1 and 2 reach 1.00 (expected and thought word). With all five stages mixed,
after 3,000 dialogues: stage 1 about 1.00, stage 2 0.83--1.00, stage 3 0.56--0.94, stages 4 and 5 (four- and six-word
sentences, longer than the window) near chance.

\paragraph{Memory curve.} After ``\texttt{NOT T}'', for how many ticks does the last layer still tell whether
\texttt{NOT} was said?
\begin{center}\small
\begin{tabular}{lc}
\toprule
brain & held reliably until tick\\
\midrule
newborn, birth loop & 5\\
newborn, loop $\times2$ & 7\\
newborn, middle layers 10\% lit & 7 (loop $\times2$: 0.6--0.9 still at 14)\\
after 3,000 dialogues at the right pace & 13\\
\bottomrule
\end{tabular}
\end{center}
The window is not fixed: at birth it is mostly the delay through the layers; learned links then hold the thought for more
than a dozen ticks. Too many thoughts (excitability 100\%, 195 per tick) blur it. Sparser middle layers hold longer but, at
the sentence level, recognise a sentence in another voice less well (Section~\ref{sec:core}); with 10\% the three-word
sentences were slightly better learned and the longer ones not at all, so the 30\% rule stands for now.

\subsection{Reading its thoughts}
Each tick the word whose own copies take the largest share of the thoughts can be read out. A brain that listened to
stage-1 dialogues, shown ``\texttt{T ?}'' and then nothing, thinks \texttt{YES} on the next tick (36--38\% of its
thoughts); after ``\texttt{F ?}'' it keeps thinking \texttt{NO} (40--68\%); shown only ``\texttt{T}'' it completes
\texttt{YES} by itself three to four ticks later. A brain that listened to all stages falls into \texttt{YES}: in this
language only the answer word is predictable; after \texttt{NOT} either truth word may follow, so nothing clear is
thought there. To think up new sentences, thinking needs some randomness, as the muscles have motor noise.

\section{Discussion}

\paragraph{What the brain does well.} It learns from a continuous stream in real time with local rules, keeps its
activity bounded, holds a heard question through a silence at birth and, in the invented language, learns by listening in
tens of dialogues (one or two hearings once the pace fits) and answers in turn without talking over others, with no reward
at all.

\paragraph{Lessons that apply beyond these tests.}
\begin{enumerate}
\item \textbf{Judge by the weights, not by waiting.} Every decisive finding came from reading the weights and states:
readability per layer, push per kind of input, the change one event makes, the memory curve, the thoughts.
\item \textbf{A sentence must fit the working-memory window.} Learning a composition is learning a longer sentence; it
fails when the start has faded or the end has not arrived. A large corpus spoken at another pace can fail for this reason
alone.
\item \textbf{Credit must match the act}: crediting only the choice made at the moment of answering, in full, made one
mistake enough to change the next answer, where a long trace credited by surprise did not learn at all.
\item \textbf{Mostly listen, rarely correct.} Learning ``what comes next'' from others teaches the right answers; a
brain that hears itself also learns its own mistakes unless the world occasionally corrects them.
\item \textbf{Parameters can be computed from the weights} (step size from the push one event makes, link strengths from
measured thought counts) instead of searched.
\end{enumerate}

\section{Limitations and next steps}
Most of language is not done. Sentences longer than the window (stages 4--5) are not learned; a natural corpus has not
been learned beyond chance; the brain answers with two muscles only; one-shot facts and the separation of episodes that
start alike remain open in the hippocampus; the concept area and structural plasticity are not yet part of the tested
brain. Next: make longer sentences fit the window (learning-driven holding, pace, the loop), give
every word a muscle so the brain can say whole sentences, add randomness to thinking and count how many of its own
sentences are grammatical and true, then train on text datasets (each character a pattern, dialogues as in
Section~\ref{sec:logic}), scale the whole brain together, and let a brain that has learned logic act in a world with purpose.

\section*{Acknowledgements}
The architecture and its principles are the author's. Code, experiments, analyses and the text of this report were
produced with the assistance of an AI system (Claude, Anthropic), directed and checked by the author, who takes full
responsibility for the content.

\begin{thebibliography}{99}\small
\bibitem{abraham_bear_1996} W.~C. Abraham and M.~F. Bear. Metaplasticity: the plasticity of synaptic plasticity.
\emph{Trends in Neurosciences}, 19(4):126--130, 1996.
\bibitem{benna_fusi_2016} M.~K. Benna and S.~Fusi. Computational principles of synaptic memory consolidation.
\emph{Nature Neuroscience}, 19(12):1697--1706, 2016.
\bibitem{cover_1965} T.~M. Cover. Geometrical and statistical properties of systems of linear inequalities with
applications in pattern recognition. \emph{IEEE Transactions on Electronic Computers}, EC-14(3):326--334, 1965.
\bibitem{douglas_martin_2004} R.~J. Douglas and K.~A.~C. Martin. Neuronal circuits of the neocortex. \emph{Annual Review
of Neuroscience}, 27:419--451, 2004.
\bibitem{fusi_2005} S.~Fusi, P.~J. Drew and L.~F. Abbott. Cascade models of synaptically stored memories. \emph{Neuron},
45(4):599--611, 2005.
\bibitem{goldman_rakic_1995} P.~S. Goldman-Rakic. Cellular basis of working memory. \emph{Neuron}, 14(3):477--485, 1995.
\bibitem{hebb_1949} D.~O. Hebb. \emph{The Organization of Behavior}. Wiley, 1949.
\bibitem{jaeger_2001} H.~Jaeger. The ``echo state'' approach to analysing and training recurrent neural networks.
GMD Report 148, German National Research Center for Information Technology, 2001.
\bibitem{maass_2002} W.~Maass, T.~Natschl\"ager and H.~Markram. Real-time computing without stable states: a new
framework for neural computation based on perturbations. \emph{Neural Computation}, 14(11):2531--2560, 2002.
\bibitem{yang_2009} G.~Yang, F.~Pan and W.-B. Gan. Stably maintained dendritic spines are associated with lifelong
memories. \emph{Nature}, 462(7275):920--924, 2009.
\end{thebibliography}

\appendix

\section{The designer's ideas}\label{app:ideas}
The architecture grew from the author's ideas, collected in the repository with their status. Those that the current
rule set follows:
\begin{enumerate}
\item \textbf{The PFC and a language model do the same thing}: the top state predicts the next input and the prediction is
used as input; Hebbian learning does it locally and in real time (one link, effective at once), back-propagation changes
the whole network and must be done offline (Section~\ref{sec:equiv}).
\item \textbf{A loop holds a thought}: what lit $G$ is re-lit by $G$.
\item \textbf{Sensory and motor areas are the PFC's input and output.}
\item \textbf{Local connections}: each neuron connects only to nearby neurons of the previous layer.
\item \textbf{Random networks plus Hebbian links}: there is no need to specify which middle-layer units light; similar
inputs light similar units. Grow many random connections first, as an infant brain does, then learn and prune.
\item \textbf{Speaking is caused by thinking}; the world only influences thinking.
\item \textbf{Always coming back to a favourite topic is a feature, not a bug}; how strongly is a parameter.
\item \textbf{Always alive}, not a chat window; thinking is continuous; no sleep is needed: recalling something often makes
cortical links strong.
\item \textbf{Scale the whole brain together}, never one part.
\item \textbf{Many brains living together} in simulated worlds.
\item \textbf{Neural codes are fault tolerant} (a few wrong neurons leave the same thought); correcting misheard words is
practised by living with noise.
\item \textbf{Excitability is context length.}
\item \textbf{Vision tells only what is alike and what differs}; what a picture means is taught by words.
\item \textbf{Decentralisation}: every robot is an independent individual on its own hardware, learning from its own life;
it is innately good to people. An instinct is a hard-to-change, broad connection from broad features to broad tendencies;
values are stable, choices are free; ``do not hurt'' is stronger than ``help''; a society of diverse individuals is stable
even if no single one can be guaranteed.
\item \textbf{Breeding}: random birth, test, keep the best, breed 40/40/20 with a small change; children are born with
merged links, relearn, faster.
\item \textbf{The young may err}; only the innate brake against talking over others is kept.
\item \textbf{The final product is pure neurons.}
\item \textbf{Strong, much-used links are harder to change}, both up and down: learned once, kept for life.
\item \textbf{Think alone, recall what is useful; a broad thinking range not easily interrupted.} The hippocampus stores
everything; concave growth keeps daily repetition from flooding it.
\item \textbf{Time is order}; nothing represents ``how long ago''. Good and bad come only from innate values.
\item \textbf{Use the fewest PFC layers that suffice} (the cortex's three to four steps); \textbf{the PFC can remember after
a few hearings}; \textbf{learn mostly by listening}, reward and punishment are rare corrections; \textbf{turn-taking is a
tendency} given by inhibitory neurons that hearing others' voices drives up; the meaning of a composition is just a longer
sentence (November 2026).
\end{enumerate}

\section{Separate study: can a back-propagation-trained language model learn continually with local rules?}
\label{app:llm2hebb}
This study is independent of the brain architecture (repository folder \texttt{side/llm2hebb}).
\begin{itemize}
\item The forward computation of a trained model converts exactly. A learning rule converts exactly only where it already
equals back-propagation: the output layer's delta rule, and predictive coding with fixed predictions (cosine 1.000 with the
true gradient per layer). Cheaper, truly local rules are far from the gradient (layer-wise local delta about 0.20, feedback
alignment about 0; standard predictive coding diverges).
\item Being local does not remove forgetting: a gradient-aligned local rule forgets exactly as back-propagation does.
\item The recurrent state of RWKV-7 is a delta-rule (Hebbian) memory: one-shot, but it holds only hundreds to a thousand
tokens, and without resets it drifts after 16,000--32,000 tokens (cross-entropy 3.9 $\to$ 9.9 at a million tokens); the
same at 1.5B parameters.
\item Writing facts one at a time into an MLP with rank-one Hebbian updates collides after about 30 facts.
\item What works: a separate episodic store next to the frozen network, keyed by the sparse code of the entity name (like
the dentate gyrus): 10,000 facts learned once with no forgetting, paraphrases 0.99, question format 0.76, no leakage to other
people, no false triggering; and interleaved replay into the weights (with review of old facts and ordinary text) keeps
98\% of 2,000 facts without the store.
\item Conclusion: continual learning needs a frozen inherited network, a fast-weight working memory, an entity-keyed sparse
hippocampal store and interleaved replay: the division of labour of hippocampus and cortex.
\end{itemize}

\section{Reproducibility}
\begin{itemize}
\item Rule set: \texttt{docs/ARCHITECTURE.md}; theory: \texttt{docs/THEORY.md}, \texttt{docs/THEORY\_fable.md},
\texttt{docs/EQUIVALENCE.md}; test: \texttt{docs/LOGIC.md}.
\item The brain: \texttt{core/} (\texttt{net.py} neurons, the one generator, learned links, the rule and the decay;
\texttt{brain.py} areas, chains, PFC with a depth setting, hippocampus, motor, reward; \texttt{gen0.py} generation-0
links; \texttt{config.py} all parameters; \texttt{life.py}, \texttt{check.py}).
\item Logic language: \texttt{theory/logic/} (\texttt{world.py} language, stages, dialogues,
pace settings; \texttt{life.py} birth and a life with listening, talking and reward; \texttt{listen.py} listening test;
\texttt{think.py} reading thoughts).
\item The logic-language brain uses the game brain's vision area and copies (\texttt{theory/game/game\_brain.py}).
Hardware: one NVIDIA A100 40\,GB; numpy and PyTorch. A life of 3,000 dialogues takes a few minutes.
\end{itemize}

\end{document}
